\scanpage{6}
\mspage{4}

Cet exemple contient le cas de $\operatorname{car}.p>0$ de l'exemple 1, si on note qu'à\deleted{\unsure{r|n}} priori, grâce à Cohen, un cristal absolu sur k,c'est pareil qu'un cristal relativement à W.

On peut aussi donner un énoncé commun aux \deleted{deux} exemples 2 et 3, en partant du cas où on se donne un $S\longrightarrow R$ non \uline{ramifié}, ce qui permet en \unsure{effett|effet} de construire encore des ``voisinages infinitésimaux'' $S_n$.

1.8. \uline{Relation avec la notion de stratification}. Les données $R,S,\underline F$ étant comme d'habitude, considérons pour chaque entier $n\ge0$ le voisinage infinitésimal $\Delta_n$ de la diagonale de $S\times_RS$, qui s'envoie dans S par les deux projections $\operatorname{pr}_1$ et $\operatorname{pr}_2$. Si E est un objet de $\underline F_S$, une n-\uline{connexion} sur E(relativement à R) est la donnée d'un isomorphisme $\operatorname{pr}_1^*(E)\simeq\operatorname{pr}_2^*(E)$ qui induit l'identité sur la diagonale. Une $\infty$-\uline{connexion} ou \uline{pseudo}-\uline{\unsure{strratification|stratification}} de E, est la donnée pour tout n d'une n-connexion, de telle façon que ces n-connexions se recollent. Enfin, une \uline{stratification} sur E est la donnée d'une pseudo-stratification qui satisfait aux conditions formelles d'une donnée de descente, quand on fait intervenir les voisinages infinitésimaux de tous ordres de \inserted{la diagonale de} $S\times_RS\times_RS$. Ces notions donnent lieu à des sorites analogues à ceux de 1.2. et 1.3. Notons que lorsque S est ``formellement non ramifié sur R'' i.e. $\underline\Omega^1_{S/R}=0$, alors toutes ces notions deviennent triviales: un objet E admet alors toujours exactement une connexion de tout ordre, donc une \unsure{seule|seube} pseudo-stratification, et celle-ci est une stratification: donc la catégorie des objets de $F_S$ \uline{munis} d'une stratification relativement à R est alors équivalente, par le foncteur ``valeur en S'', à la catégorie $\underline F_R$ elle même.(Dans tous les cas, le foncteur ``valeur en S'' est \deleted{en tous cas} fidèle).

Ceci d\inserted{éfini},on voit que pour tout cristal $\underline M$ sur S de type $\underline F$ (relativement à R), sa \deleted{\unsure{vlaeur|valeur}} valeur $\underline M(S)=M$ est un objet de $\underline F_S$ muni d'une \uline{stratification canonique} relativement à R, d'où un foncteur: cristaux relatifs de type $\underline F$ $\rightsquigarrow$ objets de $\underline F$ munis d'une stratification. La remarque de 1.4. montre d'ailleurs que ce foncteur n'est pas toujours fidèle. Il y a cependant des cas intéressants
\margin{5}
